\begingroup
\setlength{\parindent}{0pt}
\setlength{\parskip}{0.40\baselineskip}
\setlength{\abovedisplayskip}{6.5pt plus 1.5pt minus 1.5pt}
\setlength{\belowdisplayskip}{6.5pt plus 1.5pt minus 1.5pt}
\setlength{\jot}{5pt}
\clubpenalty=10000
\widowpenalty=10000

\section*{How the Proof Works}

\[
 \partial_tu+(u\cdot\nabla)u=-\nabla p+\nu\Delta u,\qquad \nabla\cdot u=0.
\]

We start with a smooth, incompressible flow of finite energy and keep its
viscosity fixed and positive. The problem is to prove that the flow remains
smooth for all time.

Strain within the fluid's flow can naturally stretch a vortex. Incompressibility
narrows the material tube as it lengthens, and sustained circulation makes the
narrowing flow spin faster. Viscosity spreads that rotation into the surrounding
fluid, so circulation can escape the contracting tube. The vorticity
\(\omega=\nabla\times u\) evolves under both effects:
\[
 \partial_t\omega+(u\cdot\nabla)\omega
 =(\omega\cdot\nabla)u+\nu\Delta\omega.
\]
The hope is simple: viscosity smooths the fluid faster than nonlinear
transport can sharpen it.

If that concentration continues, faster motion occupies an ever smaller
region. We can follow it toward vanishing scale to see whether viscosity
eventually gains an advantage. Under the Navier--Stokes rescaling,
\[
 \begin{aligned}
 u_\lambda(x,t)&=\lambda u(\lambda x,\lambda^2t),\\
 p_\lambda(x,t)&=\lambda^2p(\lambda x,\lambda^2t),
 \qquad \lambda>1.
 \end{aligned}
\]
The viscosity stays fixed. Transport and diffusion scale together, so the
equation allows the competition to persist however small the region becomes.
The velocity grows with \(\lambda\), yet the volume carrying that motion
shrinks fast enough for its energy to fall. Spatial derivatives reveal the
concentration that energy can miss. At Sobolev order \(s\), measuring \(s\)
derivatives in \(L^2\),
\[
 \begin{aligned}
 \|u_\lambda(\cdot,t)\|_{\dot H^s}
 &=\lambda^{s-1/2}\|u(\cdot,\lambda^2t)\|_{\dot H^s},\\
 E[u_\lambda](t)&=\lambda^{-1}E[u](\lambda^2t),
 \qquad E=\tfrac12\|u\|_2^2.
 \end{aligned}
\]
Below \(s=\tfrac12\), the norm falls as the flow concentrates; above it,
the norm grows. At \(s=\tfrac12\), it is invariant: the critical norm can
stay fixed while speed rises without bound and energy tends to zero.
Finite energy alone cannot exclude that behavior.

Nor does finite time exclude it. Each contraction by a fixed factor
\(\lambda\) shortens the corresponding time by \(\lambda^2\). If the
fluid could sustain successive contractions, infinitely many could fit
before a finite time \(T_*\):
\[
 \tau_j=\lambda^{-2j}\tau_0,
 \qquad
 \sum_{j=0}^{\infty}\tau_j
 =\frac{\tau_0}{1-\lambda^{-2}}<\infty.
\]
That is the possible Zeno cascade. Rescaling leaves room for it; the
regularity problem is whether the evolving fluid can actually complete it.

\needspace{14\baselineskip}
For a singularity to form in finite time, velocity would have to blow up.
That requires acceleration to blow up, then jerk, and so on, and so on.
Here these are temporal derivatives of the velocity field, with blowup
understood in spatial supremum norm as \(t\uparrow T_*\). The equation
ties each of those temporal demands to spatial derivatives:
\[
 \begin{aligned}
 \partial_tu={}&\nu\Delta u-(u\cdot\nabla)u-\nabla p,\\
 \partial_t^2u={}&\nu\Delta\partial_tu
 -(\partial_tu\cdot\nabla)u
 -(u\cdot\nabla)\partial_tu-\nabla\partial_tp,\\
 \partial_t^3u={}&\nu\Delta\partial_t^2u
 -(\partial_t^2u\cdot\nabla)u
 -2(\partial_tu\cdot\nabla)\partial_tu
 -(u\cdot\nabla)\partial_t^2u
 -\nabla\partial_t^2p.
 \end{aligned}
\]
Each further temporal order reaches another Laplacian into the preceding
response and differentiates the nonlinear terms along with it. Expanding
those responses in the original velocity exposes \(\Delta u\),
\(\Delta^2u\), \(\Delta^3u\), and higher spatial orders. The temporal
escalation is also a demand on the spatial structure of the flow.

This becomes particularly useful at the critical Sobolev order. Write
\(E_s=\tfrac12\|\Lambda^su\|_2^2\), with
\(\Lambda=(-\Delta)^{1/2}\), and \(H=E_{1/2}\). The viscous contribution
to the change of \(E_s\) is \(-2\nu E_{s+1}\). Starting from \(H\),
successive temporal derivatives expose successively higher spatial energies:
\[
 \begin{aligned}
 E_s'&=-2\nu E_{s+1}+\mathcal P_s,\\
 H'&=-2\nu E_{3/2}+\mathcal P_{1/2},\\
 H''&=4\nu^2E_{5/2}-2\nu\mathcal P_{3/2}
       +\partial_t\mathcal P_{1/2},
 \end{aligned}
\]
where the nonlinear contribution is
\[
 \mathcal P_s=-\operatorname{Re}\langle\Lambda^su,
 \Lambda^s((u\cdot\nabla)u+\nabla p)\rangle_{L^2}.
\]
The pattern continues through \(E_{7/2},E_{9/2},\ldots\). The nonlinear
terms remain coupled to those rising spatial orders throughout the evolution.

Those rising spatial orders are also the orders that give smoothness.
For any finite number of continuous derivatives, a sufficiently high
Sobolev order is enough:
\[
 H^m(\mathbb R^3)\hookrightarrow C^k(\mathbb R^3)
 \quad\text{when }m>k+\tfrac32,
\]
and finite energy supplies the low frequencies. Taking all the orders
together gives
\[
 L^2\cap\bigcap_{n\ge0}\dot H^{n+1/2}
 =\bigcap_{m\ge0}H^m
 =H^\infty\hookrightarrow C^\infty.
\]
The spatial requirements of smoothness have emerged from following the
temporal requirements of a singularity. If the equation carries this hierarchy through
\(T_*\), the velocity reaches that time with enough regularity to continue.
The question has become whether the derivatives generated by the initial
flow can remain together in this intersection at the endpoint.

If a singularity were to form, we would expect smoothness at every earlier
instant but loss of a fixed Sobolev norm on the approach to \(T_*\).
Velocity and its temporal responses would become unbounded even while
energy remained finite. All the derivatives could exist before \(T_*\)
without reaching it as derivatives of a smooth limiting velocity.

\needspace{11\baselineskip}
We begin the construction with the derivatives determined by the initial
velocity \(u_0\in H^\infty_\sigma(\mathbb R^3)\). Incompressibility fixes
pressure through \(-\Delta p=\partial_i\partial_j(u_i u_j)\), with its
decay normalization. The differentiated equations above already show how
one temporal row generates the next. Writing
\(V_n=\partial_t^nu\) and \(Q_n=\partial_t^np\) retains the full pattern:
\[
 \begin{aligned}
 V_{n+1}&=\nu\Delta V_n
 -\sum_{a=0}^n\binom na(V_a\cdot\nabla)V_{n-a}-\nabla Q_n,\\
 Q_n&=R_iR_j\sum_{a=0}^n\binom na V_{a,i}V_{n-a,j}.
 \end{aligned}
\]
Here \(R_i\) are the Riesz transforms, with repeated spatial indices
summed. At \(t=0\), the recurrence starts from \(V_0=u_0\); spatial
differentiation fills out each temporal row. This gives a mixed velocity
and pressure jet whose finite pieces agree wherever they overlap.

Fixing a temporal depth \(N\) and spatial height \(M\) selects a finite
rectangle \(\mathcal R_{N,M}[u_0;T_*]\). The datum-generated rectangle
theorem, Theorem~\ref{thm:datum-rectangles}, gives adjacent rows square
integrable in \(H^{M+1}\) and \(H^{M-1}\) all the way to \(T_*\).
These finite pieces agree
on overlaps and assemble into the compatible intersection:
\[
 \begin{aligned}
 \mathcal I[u_0]
 &:=\bigl(\mathcal R_{N,M}[u_0;T_*]\bigr)_{N,M\ge0},\\
 u&\in\bigcap_{r,m\ge0}H^r(0,T_*;H^m(\mathbb R^3)).
 \end{aligned}
\]
The intersection's adjacent temporal coordinates give
\[
 \begin{aligned}
 V_n&\in L^2(0,T_*;H^{M+1}),\\
 \partial_tV_n=V_{n+1}&\in L^2(0,T_*;H^{M-1}),
 \end{aligned}
 \qquad
 V_n\in C([0,T_*];H^M).
\]
The Hilbert-triple trace theorem gives the last inclusion. Compatibility
identifies the traces across all rectangles, leaving one endpoint in every
spatial Sobolev space:
\[
 \begin{aligned}
 u_*&=\lim_{t\uparrow T_*}u(t)\quad\text{in every }H^m,\\
 u_*&\in H^\infty_\sigma\hookrightarrow C^\infty.
 \end{aligned}
\]
Local existence starts a smooth flow from \(u_*\), and uniqueness joins
it to the original solution. The supposed maximal lifespan extends
beyond \(T_*\). A finite \(T_*\) is impossible.

We can now extract the finite estimates and use them. At any finite
derivative orders and any finite time \(T\), Sobolev embedding gives
\[
 \partial_t^ru\in C([0,T];H^m)
 \hookrightarrow L^\infty(0,T;W^{k,p}),
 \qquad m>k+\tfrac32-\tfrac3p,
\]
for \(r,k\ge0\) and \(2\le p\le\infty\). In particular, the total
strain accumulated in finite time is finite:
\[
 \int_0^T\|\nabla u(t)\|_\infty\,dt
 \le T\sup_{0\le t\le T}\|\nabla u(t)\|_\infty<\infty.
\]
Returning to a locally circular material tube, axial strain
\(\sigma\) gives \(\dot r/r=-\sigma/2\), with
\(\sigma\le\|\nabla u\|_\infty\). Its radius obeys
\[
 r(t)\ge r(0)\exp\!\left(
 -\frac12\int_0^t\|\nabla u(s)\|_\infty\,ds
 \right)>0.
\]
Only a finite amount of stretching can accumulate in a finite time.
The tube cannot be squeezed to zero radius there.

\endgroup
